\documentclass[12pt]{amsart} 

\input{./latex-preamble/cm-preamble.tex}

%\graphicspath{ {./diagrams/} }

\newtheorem*{lemma}{Lemma}

\begin{document}

{Hartshorne} \hfill {2-7-2025}

\bigskip\bigskip

\begin{center}

{\sc Chapter II.2-3, Week 3}

\end{center}

\begin{center}

  {Noah Parker} %\footnote{Collaborated with Kalev Martinson and Frank :).}}
    
\end{center}

\bigskip

\begin{enumerate}
	\item[2.14.] {
      \begin{enumerate}
        \item Let $S$ be a graded ring.
          Show that $\Proj S=\0$ if and only if every element of $S_+$ is nilpotent.
        \item Let $\vphi:S\to T$ be a graded homomorphism of graded rings (preserving degrees).
          Let $U=\{\p\in \Proj T~:~\p\not\supset\vphi(S_+)\}$.
          Show that $U$ is an open subset of $\Proj T$, and show that $\vphi$ determines a natural morphism $f:U\to \Proj S$.
        \item The morphism $f$ can be an isomorphism even when $\vphi$ is not.
          For example, suppose that $\vphi:S_d\to T_d$ is an isomorphism for all $d\geqs d_0$, where $d_0$ is any integer.
          Then show that $U=\Proj T$ and the morphism $f:\Proj T\to \Proj S$ is an isomorphism.
      \end{enumerate}
    }
\end{enumerate}

\begin{proof}
  (a) Denote $\nil' S := \cap_\p\p$, where $\p$ ranges over all homogeneous prime ideals of $S$.
  It suffices to show that $\nil' S=\nil S$.
  % First, we show that homogeneous prime ideals are prime in the usual sense.

  % Fix a homogeneous prime ideal $\p$.
  % Suppose that $ab\in \p$, and take $a=\sum_{i=1}^ra_i$, $b=\sum_{j=1}^sb_j$.
  % Then 
  % \[
  %   ab = \sum_{\substack{1\leqs i\leqs r \\ 1\leqs j\leqs s}}a_ib_j\in\p
  % \] 
  % gives us $a_rb_s\in \p$, so that $a_r\in\p$ or $b_s\in \p$.
  % Without loss of generality, $a_r\in \p$, implying $(a-a_r)b = ab-a_rb\in \p$.
  % Induction allows us to reduce to the $a,b$ homogeneous of degree 0 case, whence $\p$ is prime.
  % It follows that $\nil S\subset \nil' S$.

  % Suppose $f\in \nil S$.
  % Write $f=\sum_{i=1}^d f_i$, and take $n>0$ with $f^n=0$.
  % Then $f_d^n=0$ gives us that $f_d\in \nil S$.
  % Replacing $f$ by $f-f_d\in \nil S$ and repeating the process gives us that each $f_i\in \nil S$.
  % Thus, $\nil S$ is a homogeneous prime ideal.
  For any prime ideal $\p$, the homogenization $h(\p)\subset \p$ generated by its homogeneous elements is clearly a homogeneous prime ideal.
  Then we have
  \[
    \nil' S \supset \nil S=\bigcap_{\p}\p \supset \bigcap_{\p}h(\p)=\nil' S,
  \] 
  %since each homogeneous prime ideal is prime in the usual sense.
  which gives us $\nil'S\subset \nil S$.
  We conclude that $\nil S=\nil' S$.
  Now, $\Proj S=\0$ iff $S_+\subset \nil' S=\nil S$, which is exactly our result.

  (b) Let $\p\in \Proj T$ be a homogeneous prime ideal.
  Then $\vphi^{-1}(\p)$ is a homogeneous prime ideal, and $\vphi^{-1}(\p)\supset S_+$ implies 
  \[
    \p\supset \vphi(\vphi^{-1}(\p))\supset \vphi (S_+),
  \]
  i.e. $\p\notin U$.
  Conversely, $\p\supset \vphi(S_+)$ implies 
  \[
    \vphi^{-1}(\p)\supset \vphi^{-1}(\vphi(S_+))\supset S_+,
  \]
  so that $\vphi^{-1}(\p)\notin \Proj S$.
  We conclude that $U$ is the inverse image of $\Proj S$ under the mapping of homogeneous prime ideals $\p\mapsto \vphi^{-1}$.
  Hence, we have a natural map $f:U\to\Proj S:\p\mapsto \vphi^{-1}(\p)$.

  Fix a homogeneous $g\in T_+$.
  Then
  \begin{align*}
    \vphi^{-1}(\p)\ni g\implies \p\supset \vphi(\vphi^{-1}(\p))\ni \vphi(g) \implies \vphi^{-1}(\p)\supset \vphi^{-1}(\vphi(g))\ni g
  \end{align*}
  gives us that
  \begin{align*}
    f^{-1}(D_+(g)) &= D_+(\vphi(g))\subset U.
  \end{align*}
  Hence, $f:U\to \Proj S$ is continuous.
  We must now define the map $f^\sharp:\O_S\to f_*(\O_T|_U)$ of sheaves.

  It suffices to define $f^\sharp$ on basic open sets.
  Fix a homogeneous $g\in S_+$, and note that 
  \[
    f_*(\O_T|_U)(D_+(g)) = \O_T(f^{-1}(D_+(g))\cap U) = \O_T(D_+(\vphi(g)))\cong T_{(\vphi(g))}.
  \]
  Recall as that $\vphi$ is degree preserving, so each of its localizations are as well.
  In particular, $\vphi_g:S_g\to T_g$ restricts to a map $\vphi_g:S_{(g)}\to T_{(\vphi(g))}$.
  Then the following commutative diagram gives us our sheaf map on $D_+(g)$.
  \[
    \begin{tikzcd}[column sep = 1.2cm]
      \O_S(D_+(g))\ar[d, "\sim" {rotate=90, anchor=north}]\ar[r,"f^\sharp(D_+(g))"]
      &f_*(\O_T|_U)(D_+(g))\ar[d, "\sim" {rotate=90, anchor=north}] \\
      S_{(g)}\ar[r,"\vphi_g"']&T_{(\vphi(g))}
    \end{tikzcd}
  \] 
  To show that these maps glue together, we note that $D(g)\cap D(h)= D(gh)$, and consider the following commutative diagram.
  \[
    \begin{tikzcd}[column sep = 0.05cm,row sep = 0.4cm]
      S_{(g)}\ar[rr,"\vphi_g"]\ar[dr]
        &&S_{(\vphi(g))}\ar[dr]\\
      &S_{(gh)}\ar[rr,"\vphi_{gh}"] 
        &&S_{(\vphi(gh))} \\
      S_{(h)}\ar[ur]\ar[rr,"\vphi_h"']
        &&S_{(\vphi(h))}\ar[ur]
    \end{tikzcd}
  \] 
  This defines our sheaf homomorphism on all of $\Proj S$, so that we're done.

  % Let $s\in \O_S(V)$, and define 
  % \[
  %   f^\sharp(V)(s)(\p) = \vphi_{f(\p)}(s(f(\p)))
  % \]
  % for all $\p\in f^{-1}(U\cap V)\subset \Proj T$.
  % $\vphi$ is degree preserving, so each of its homogeneous localizations $\vphi_{f(\p)}$ is degree preserving.

  % If $s|_W(\q)=a/g\in (\O_S)_{(\q)}$ for $a,g\in S$ homogeneous of the same degree and $g\notin \q$ for all $\q\in U\cap W$, then 
  % \[
  %   f^\sharp(W)(s|_W)(\p) = \vphi_{f(\p)}(a/g)\in \vphi_{f(\p)}(\O_S)_{(f(\p))}\subset (\O_T)_{(\p)}
  % \]
  % for all $\p\in f^{-1}(U\cap W)$.
  % $\vphi_{f(\p)}(a/g)=\vphi(a)/\vphi(g)$, with $\vphi(a),\vphi(b)\in T$ homogeneous of the same degree since $\vphi$ is degree preserving, and $\vphi(g)\notin \vphi(f(\p))\subset \p$.
  % Each $f^\sharp_\p=\vphi_\p:(\O_S)_{(f(\p))}\to (\O_T)_{(\p)}$ is a local homomorphism, since $\vphi^{-1}(\p)=f(\p)$ by definition.
  % Naturality of $f^\sharp$ with respect to restriction is trivial, so we conclude that $(f,f^\sharp)$ is indeed a morphism of locally ringed spaces.

  (c) Fix $d_0$ such that $\vphi:S_d\to T_d$ is an isomorphism for all $d\geqs d_0$.
  We wish to show that $U=\Proj T$, so it suffices to show that $\sqrt{\vphi(S_+)}=T_+$.
  One inclusion is clear since $\vphi$ is degree preserving.
  Surjectivity of $\vphi:S_d\to T_d$ gives us $\vphi(S_+)\supset \oplus_{d\geqs d_0}T_d$.
  Then, for any element $a\in T_+$, $a^{d_0}$ is of degree at least $d_0$, so $a^{d_0}\in \vphi(S_+)$ implies $a\in \sqrt{\vphi(S_+)}$.
  We conclude that $\sqrt{\vphi(S_+)}=T_+$, so $U=\Proj T$.

  % It suffices to show that each map of stalks $\vphi_{f(\p)}:S_{(f(\p))}\to T_{(\p)}$ is an isomorphism.
  % Recall that $\vphi_{f(\p)}(a/g)=\vphi(a)/\vphi(g)$, whenever $a,g\in S$ are homogeneous of the same degree and $g\notin f(\p)$.
  % $f(\p)\in \Proj S$ gives us $h\in S_+\setminus f(\p)$.
  % Replacing $a/g$ by $ah^{d_0}/gh^{d_0}$, we may assume that $a,g$ are of degree $\geqs d_0$.
  % Then
  % \begin{align*}
  %   \vphi_{f(\p)}(a/g) &= \vphi_{f(\p)}(a'/g')\\
  %   \implies ~\frac{\vphi(a)}{\vphi(g)}
  %                      &=\frac{\vphi(a')}{\vphi(g')}.
  % \end{align*}
  % Thus, we can find $t\in T_+\setminus \p$ such that
  % \[
  %   t(\vphi(a)\vphi(g')-\vphi(a')\vphi(g))=0.
  % \] 
  % Without loss of generality, we may assume $t\in T_d$ for some $d\geqs d_0$.
  % Then $T_d=\vphi(S_d)$ gives us $s\in S_d$ with $t=\vphi(s)$, so
  % \begin{align*}
  %   \vphi(t(ag'-a'g))=0\implies \frac{a}g=\frac{a'}{g'}
  % \end{align*}
  % by injectivity of $\vphi:\oplus_{d\geqs d_0}S_d\to \oplus_{d_\geqs d_0}T_d$.
  % Thus, $\vphi_{f(\p)}:S_{(f(\p))}\to T_{(\p)}$ is injective.

  % Take $b/h\in T_{(\p)}$. 
  % By the above argument, we may assume $b,h$ are of degree $d\geqs d_0$.
  % Then we can take homogeneous $a,g$ of degree $d\geqs d_0$ such that $\vphi(a)=b$ and $\vphi(g)=h$, so that
  % \begin{align*}
  %   \vphi_{f(\p)}(a/g)= \frac{\vphi(a)}{\vphi(g)} = \frac{b}h
  % \end{align*}
  % gives surjectivity of $\vphi_{f(\p)}$.
  % We conclude that $\vphi_{f(\p)}$ is an isomorphism for each $\p\in \Proj T$, so that $\vphi$ is an isomorphism of schemes.

  It suffices to show that $f^\sharp$ is an isomorphism on open sets $D(g)$.
  Fix homogeneous $g\in S_+$.
  Since $D(g^n)=D(g)$ for all $n>0$, we may assume $g$ is of degree $\geqs d_0$.
  We show the map $\vphi_g:S_{(g)}\to T_{(\vphi(g))}$ is surjective.
  For any $b/\vphi(g)^n\in (B_A)_\p$, $b\vphi(g)=\vphi(a)$ with $a\in A$ gives us 
  \[
    \vphi_g(a/g^{n+1})=\vphi(ag)/\vphi(g^{n+1})=b/\vphi(g)^n\in \im\vphi_g.
  \]
  It suffices to check injectivity.
  If $\vphi(a/g)=0$, then we can find $n\geqs 1$ such that $\vphi(g)^n\vphi(a)=0$.
  Then $\vphi(g^na)=0$ implies $g^na=0$ by injectivity of $\vphi$, whence $a/g =0$ in $S_{(g)}$.
  The result follows.
\end{proof}


% \begin{enumerate}
% 	\item[2.16.] {
%       $(*)$ Let $X$ be a scheme, $f\in \Gamma(X,\O_X)$, and define $X_f:=\{x\in X:f_x\notin m_x\subset \O_x\}$.
%       \begin{enumerate}
%         \item If $U=\Spec B$ is an open affine subscheme of $X$, and if $\ol{f}\in B=\Gamma(U,\O_X|_U)$ is the restriction of $f$, show that $U\cap X_f=D(\ol{f})$.
%           Conclude that $\ol{f}$ is an open subset of $X$.
%         \item Assume that $X$ is quasi-compact.
%           Let $A=\Gamma(X,\O_X)$, and let $a\in A$ be an element whose restriction to $X_f$ is 0.
%           Show that for some $n>0$, $f^na=0$.
%         \item Now assume that $X$ has a finite cover by open affines $U_i$ such that each intersection $U_i\cap U_j$ is quasi-compact. (This hypothesis is satisfied, for example, if $\Sp(X)$ is noetherian.)
%           Let $b\in \Gamma(X_f,\O_{X_f})$.
%           Show that, for some $n>0$, $f^nb$ is the restriction of an element of $A$.
%         \item With the hypotheseis of (c), conclude that $\Gamma(X_f,\O_{X_f})\cong A_f$.
%       \end{enumerate}
%     }
% \end{enumerate}

\begin{enumerate}
	\item[2.18.] {
      In this exercise, we compare some properties of a ring homomorphism to the induced morphism of the spectra of the rings.
      \begin{enumerate}
        \item Let $A$ be a ring, $X=\Spec A$, and $f\in A$.
          Show that $f$ is nilpotent if and only if $D(f)$ is empty.
        \item Let $\vphi:A\to B$ be a homomorphism of rings, and let $f:Y=\Spec B\to X=\Spec A$ be the induced morphism of affine schemes.
          Show that $\vphi$ is injective if and only if the map of sheaves $f^\sharp:\O_X\to f_*\O_Y$ is injective.
          Show furthermore in that case $f$ is \emph{dominant}, i.e., $f(Y)$ is dense in $X$.
        \item With the same notation, show that if $\vphi$ is surjective, then $f$ is a homeomorphism of $Y$ onto a closed subset of $X$, and $f^\sharp:\O_X\to f_*\O_Y$ is surjective.
        \item Prove the converse to (c).
          Namely, if $f:Y\to X$ is a homeomorphism onto a closed subset, and $f^\sharp:\O_X\to f_*\O_Y$ is surjective, then $\vphi$ is surjective.
      \end{enumerate}
    }
\end{enumerate}

\begin{proof}
  (a) The result is trivial, as $D(f)=\0$ iff $f\in \cap_\p\p=\nil f$.

  (b) In the language of coherent sheaves, we have $\O_X = \wtilde{A}$ and $f_*\O_Y=f_*\wtilde{B} = \wtilde{B_A}$.
  Then $f^\sharp_\p:\O_{X,\p}\to (f_*\O_Y)_\p$ corresponds to the map $\vphi_\p:A_\p\to (B_A)_\p$.
  By CA, we have that $\vphi$ is injective iff each $\vphi_\p$ is injective, which by the above occurs iff $f^\sharp:\O_X\to f_*\O_Y$ is injective.

  It suffices to show that $f(Y)$ is dense in $X$.
  Fix $g\in A$.
  Then
  \begin{align*}
    f(Y)\cap D(g)=\0 &\iff f^{-1}(D(g))=\0 \\
                     &\iff D(\vphi(g))=\0 \\
                     &\iff \vphi(g)\in \nil B \\
                     &\iff g\in \nil A \\
                     &\iff D(g)=\0
  \end{align*}
  by (a) and injectivity of $\vphi$ respectively.
  Any nonempty open $U\subset X$ contains a nonempty distinguished open of $X$, so $f(Y)$ is dense in $X$.
  
  % Clearly $f^\sharp$ injective implies $\vphi=f^\sharp(X)$ is injective.
  % Suppose conversely that $\vphi$ is injective.
  % It suffices to show that each map $f^\sharp_{f(\p)}=\vphi_{f(\p)}$ of stalks is injective, for $\p\in Y$.
  % 
  % % Let $a/g\in \ker\vphi_{f(\p)}$, so there is $t\in B\setminus \p$ with $t\vphi(a)=0$.
  % % If $t=\vphi(s)\in \vphi(A\setminus f(\p))\subset B\setminus \p$, then $0=t\vphi(a)=\vphi(sa)$ implies $sa=0$ in $A$, whence $a/g=0$ in $A_{f(\p)}$.
  % % Suppose instead that $\vphi(s)\vphi(a)\neq 0$ for all $s\in A\setminus f(\p)$.

  % We give $B$ the structure of an $A$-module by $a\cdot b = \vphi(a)b$.
  % Then $a\cdot \vphi(a')=\vphi(a)\vphi(a')=\vphi(aa')$ gives us that $\vphi$ is an injective morphism of $A$-module.
  % Flatness of the $A$-module $A_{f(\p)}$ gives us that $\vphi_{f(\p)}:A_{f(\p)}\hookrightarrow A_{f(\p)}\otimes_A B$ is injective.
  % We have an isomorphism $A_p\otimes_A B\xrightarrow{\sim}S^{-1}B$, with $S=\vphi(A\setminus f(\p))\subset B\setminus \p$.
  % Then we have an injection $S^{-1}B\hookrightarrow (S^{-1}B)_{S^{-1}\p}\cong B_\p$.
  % Composing these together, we see that
  % \[
  %   \vphi_{f(\p)}:
  %   A_{f(\p)}\hookrightarrow 
  %   A_{f(\p)}\otimes B \xrightarrow{\sim}
  %   S^{-1}B \hookrightarrow
  %   (S^{-1}B)_{S^{-1}\p}\xrightarrow{\sim}
  %   B_\p
  % \] 
  % is an injection.

  % It suffices to show that, in this case $f(Y)$ is dense in $X$.
  % Suppose $V(\a)\supset f(Y)$ is closed, and write $V(\a)^c = \cup_{g\in \a}D(g)$.
  % Fix $g\in \a$, and note that $D(g)\cap f(Y)=\0$.
  % Hence, $g\in f(\p)=\vphi^{-1}(\p)$ for all $\p\in X$.
  % Thus, $\vphi(g)\in\vphi(\vphi^{-1}(\p))\subset \p$ for all $\p\in X$ implies $\vphi(g)$ is nilpotent.
  % Taking $n>0$ with $\vphi(g)^n=0$, we see that $\vphi(g^n)=0$, whence $g^n=0$ by injectivity of $\vphi$.
  % We conclude that $g$ is nilpotent, so that $D(g)=\0$.
  % Hence, $V(\a)= X$, showing that $f(Y)^-=X$ is dense.

  (c) First, suppose that $\vphi:A\to B$ is surjective.
  Localization is exact, so each of the stalks $\vphi_\p:A_\p\to (B_A)_\p$ is surjective, whence $f^\sharp:\O_X\to f_*\O_Y$ is surjective.
  Surjectivity of $\vphi$ gives us
  \[
    \p\supset \a\implies\vphi^{-1}(\p)\supset \vphi^{-1}(\a)\implies \p=\vphi(\vphi^{-1}(\p))\supset \vphi(\vphi^{-1}(\a))=\a,
  \]
  so that 
  \begin{align*}
    f(V(\a)) &= \{\vphi^{-1}(\p)\in X ~:~ \p\supset \a\} \\
             &= f(Y)\cap V(\vphi^{-1}(\a)).
  \end{align*}
  In particular,
  \[
    f(Y)=f(Y)\cap V(\ker\vphi)\subset V(\ker\vphi).
  \] 

  Take $\q\in V(\ker\vphi)$, and suppose there exists $a\in \vphi^{-1}(\vphi(\q))\setminus \q$.
  Then there is some $a'\in \q$ with $\vphi(a)= \vphi(a')$, so that $a-a'\in \ker\vphi\setminus \q=\0$, a contradiction.
  Thus, $\vphi^{-1}(\vphi(\q))=\q$.

  We now show that $\vphi(\q)$ is prime.
  Suppose $bb'\in \vphi(\q)$, and take $a,a'\in A$ such that $\vphi(a)=b$ and $\vphi(a')=b'$.
  Then $aa'\in \vphi^{-1}(\vphi(\q))=\q$ implies $a\in \q$ or $a'\in \q$.
  Hence, $b\in \vphi(\q)$ or $b'\in \vphi(\q)$ gives us that $\vphi(\q)$ is prime, so that $g:V(\ker \vphi)\to Y:\q\to\vphi(\q)$ is a set-theoretic inverse to $f:Y\to V(\ker\vphi)\subset X$.
  In particular, $f(Y)=V(\ker\vphi)$.

  It suffices to show that $g:V(\ker\vphi)\to Y$ is continuous.
  We write
  \begin{align*}
    f(V(\a)) = f(V)\cap V(\vphi^{-1}(\a)) = V(\ker\vphi + \vphi^{-1}(\a)),
  \end{align*}
  so that $g^{-1}(V(\a))= V(\ker\vphi + \vphi^{-1}(\a))$ is closed in $X$.
  We conclude that $f$ is a homeomorphism onto its closed image $f(Y)=V(\ker\vphi)$.

  (d) Suppose now that $f:Y\to X$ is a homeomorphism onto a closed subset, and $f^\sharp:\O_X\to\O_Y$ is surjective.
  Let $Z=\Spec(A/\ker \vphi)$, so that the universal property of quotients supplies $\ol{\vphi}:A/\ker\vphi\to B$ such that $\vphi = \ol{\vphi}\circ \pi$.
  We denote the associated morphism of sheaves $\ol{f}:Y\to Z$.
  Note that $\pi:A\to A/\ker\vphi$ is surjective, so (c) gives us that the associated map $g:Z\to V(\ker\vphi)\subset X$ is a homeomorphism.

  $\ol{\vphi}$ is injective, so (b) implies $\ol{f}^\sharp:\O_Z\to \O_Y$ is injective, and $\ol{f}(Y)=f(Y)$ is dense.
  $f(Y)=f(Y)^-$ is closed, so $\ol{f}:Y\to Z$ is surjective.
  $f = g\circ \ol{f}$ with $f,g$ homeomorphisms onto their image implies $\ol{f}$ is a homeomorphism onto its image $\ol{f}(Y)=Z$.
  Additionally, $f^\sharp = \ol{f}^\sharp g^\sharp$ surjective implies $\ol{f}^\sharp$ is surjective, so that $\ol{f}^\sharp:\O_Z\to \O_Y$ is an isomorphism of sheaves.
  Then 
  \[
    \vphi =\ol{\vphi}\circ \pi = \ol{f}^\sharp(Z)\circ \pi
  \]
  is the composition of a surjective map with an isomorphism, and is therefore surjective itself.
\end{proof}

\begin{enumerate}
	\item[3.1.] {
      Show that a morphism $f:X\to Y$ is locally of finite type if and only if \emph{every} affine subset $V=\Spec B$ of $Y$, $f^{-1}(V)$ can be covered by open affine subsets $U_j=\Spec A_j$, where each $A_j$ is a finitely generated $B$-algebra.
    }
\end{enumerate}

\begin{proof}
  One direction is clear.
  Suppose conversely that $f:X\to Y$ is locally of finite type, and fix $\Spec A\cong U\subset X$ and $\Spec B\cong V\subset Y$ affine, with $f(U)\subset V$.
  To proceed, we must first prove the following lemma.

  \begin{lemma}
    Fix a sheaf morphism $f:X\to Y$ of finite type, $\Spec A\cong U\subset X$, $\Spec B\cong V\subset Y$, $f(U)\subset Y$ as above, and denote the map of rings associated to $f|_U:U\to V$ by $\vphi:B\to A$.
    There exist $f_{11},\ldots, f_{rs}\in A$ and $g_1,\ldots, g_r\in B$ such that
    \begin{align*}
      \Spec A &= D(f_{11})\cup\cdots \cup D(f_{rs}), \\
      \Spec B &= D(g_1)\cup \cdots \cup D(g_r),
    \end{align*}
    \[
      f(D(f_{ij}))\subset D(g_i),
    \] 
    and the induced map
    \begin{align*}
      \vphi_{g_i}: B_{g_i}\to A_{f_{ij}}
    \end{align*}
    is of finite type.
  \end{lemma}

  \begin{proof}
    We have an affine cover $\{V_i=\Spec B_i\}$ of $X$ such that each $f^{-1}(V_i)$ can be covered by open affines $U_{ij}=\Spec A_{ij}$, where $A_{ij}$ is a finitely generated $B_i$ algebra.
    Fix some $i$ and write
    \begin{align*}
      V_i\cap V &= \bigcup_{D(g_\beta)\subset V_i\cap V}D(g_\beta),
    \end{align*}
    where $D(g_\beta)$ is affine in both $V_i$ and $V$.
    Then $f^{-1}(D(g_\beta))=D(\vphi(g_\beta))$ and fixing some $j$ allows us to write
    \begin{align*}
      U_{ij}\cap U &= \bigcup_{D(f_\alpha)\subset U_{ij}\cap U}D(f_\alpha) \\
             &= \bigcup_{\substack{D(f_\alpha) \subset U_{ij}\cap U \\ D(g_\beta)\subset V_i\cap V}} 
             D(f_\alpha\vphi(g_\beta)),
    \end{align*}
    where each $D(f_\alpha)$ and $D(\vphi(g_\beta))$ is distinguisheed in both $U$ and $U_{ij}$.
    We denote $f_{\alpha\beta}:= f_{\alpha} \vphi(g_\beta)$.
    Then $f|_{D(f_{\alpha\beta})}:D(f_{\alpha\beta}) \to D(g_\beta)$ induces the map $\vphi_{g_\beta}: B_{g_\beta}\to A_{f_{\alpha\beta}}$.

    For any $g_\beta\in B$, $f_\alpha\in A$, and $f_{\alpha\beta}\in A$ as above, we write $g'_\beta\in B_i$, $f'_\alpha\in A_{ij}$, and $f'_{\alpha\beta}\in A_{ij}$ for the corresponding elements which determine the same distinguished subset.
    Then
    \begin{align*}
      (B_i)_{g'_\beta} 
      = \O_{V_i}(D(g'_\beta)) 
      = \O_V(D(g_\beta))
      = B_{g_\beta},
    \end{align*}
    \begin{align*}
      (A_{ij})_{f'_\alpha} 
      = \O_{U_{ij}}(D(f'_\alpha))  
      = \O_U(D(f_\alpha)) 
      = A_{f_\alpha},
    \end{align*}
    and
    \begin{align*}
      (A_{ij})_{f'_{\alpha\beta}}
      = A_{f_{\alpha\beta}}
      = A_{f_\alpha}\otimes_B B_{g_{\beta}}.
    \end{align*}
    Then $A_{ij}$ a finitely generated $B_i$-algebra and $A_{f_\alpha}$ a cyclic $A_{ij}$-algebra gives that $A_{f_\alpha}$ is a finitely generated $B_i$-algebra.
    That is, there exists a surjective map
    \begin{align*}
      B_i[x_1,\ldots, x_{n_\alpha}]~&\rightarrow A_{f_\alpha}.
    \end{align*}
    Right exactness of $-\otimes_B B_{g_\beta}$ gives us a surjective map
    \begin{align*}
      B_{g_\beta}[x_1,\ldots, x_{n_\alpha}]~&\rightarrow A_{f_{\alpha\beta}},
    \end{align*}
    whence $A_{f_{\alpha\beta}}$ is a finitely generated $B_{g_\beta}$-algebra.

    Combining these covers of the $V_i\cap V$ and $U_{ij}\cap U$ gives us a cover $\{D(g_\beta)\}_\beta$ of $\Spec B$, and covers $\{D(f_{\alpha\beta}) \}_\alpha$ of each $f^{-1}(D(g_\beta))=D(\vphi(g_\beta))$, where $\vphi_{g_\beta}:B_{g_\beta}\to A_{f_{\alpha\beta}}$ is of finite type.
    Since $\Spec B$ and $D(\vphi(g_\beta))$ are affine, we may assume each of these covers is finite, completing our result.
  \end{proof}

  Thus, it suffices to show the following lemma.

  \begin{lemma}
    Let $\vphi:B\to A$ be a morphism of rings such that $\vphi_i:B\to A_{f_i}$ is of finite type for each $1\leqs i\leqs r$, where $(f_1,\ldots, f_r)=(1)$.
    Then $\vphi$ is of finite type.
  \end{lemma}

  \begin{proof}
    Take $\{a_{ij}=a'_{i j}/f_i^{e_{ij}}\}_{j=1}^{r_i}$ to be a set of generators for $A_{f_i}$ as a $B$-algebra. %, so that $\{a_{ij}'\}_{j=1}^{s_i}\cup\{f_i^{-1}\}$ is another set of generators for
    Without loss of generality, $r_i=r$ for all $1\leqs i\leqs r$.
    Define
    \[
      A' := \vphi(B)[a'_{ij}~:~1\leqslant i,j\leqs r],
    \] 
    noting that
    \begin{align*}
      A'_{f_i} &= \vphi(B)[f_i^{-1},a'_{ij}~:~\leqslant i,j\leqs r]= A_{f_i}.
    \end{align*}
    Now, consider the quotient $A/A'$.
    It suffices to show that $A/A'=0$.
    For any prime ideal $\p\in \Spec A$, we can choose $i$ such that $f_i\notin \p$.
    Then
    \begin{align*}
      (A/A')_\p &= \frac{A_\p}{A'_\p} \\
                &= \frac{(A_{f_i})_\p} {(A'_{f_i})_\p} \\
                &= \left(\frac{A_{f_i}}{A_{f_i}}\right)_\p \\
                &= 0_\p \\
                &= 0
    \end{align*}
    gives us that in fact $A/A'=0$, since being zero is a local condition.
    We conclude that $A$ is a finitely generated $B$-algebra, as desired.
  \end{proof}
\end{proof}
 
 \begin{enumerate}
 	\item[3.2.] {
       A morphism $f:X\to Y$ of schemes is \emph{quasi-compact} if there is a cover of $Y$ by open affines $V_i$ such that $f^{-1}(V_i)$ is quasi-compact for each $i$.
       Show that $f$ is quasi compact if and only if for \emph{every} open affine subset $V\subset Y$, $f^{-1}(V)$ is quasi-compact.
     }
 \end{enumerate}
 
 \begin{proof}
   One implication is clear.
   Conversely, suppose there exists a cover $\{V_i=\Spec B_i\}$ of $Y$ by open affines such that $f^{-1}(V_i)$ is quasi-compact for each $i$, and take an open affine $V=\Spec B\subset Y$.
   Fix $i$ and write 
   \begin{align*}
     V\cap V_i&= \bigcup_j D(g_{ij}),
   \end{align*}
   where we range over $D(g_{ij})\subset V\cap V_i$ which are affine in both $V$ and $V_i$.
   Then
   \begin{align*}
     V&= \bigcup_{i,j}D(g_{ij}) = \bigcup_{1\leqs i,j\leqs n}D(g_{ij})
   \end{align*}
   by quasi-compactness of $V_i$.
   A finite union of compact sets is compact, so it suffices to show the following lemma.%that $f^{-1}(D(g))$ is compact for each $g\in B_i$.

   \begin{lemma}
     Let $f:X\to Y$ be a map of schemes, and $\Spec A\cong U\subset Y$ an affine open such that $f^{-1}(U)$ is quasi-compact.
     For any $h\in A$, $f^{-1}(D(h))$ is also quasi-compact.
   \end{lemma}

   \begin{proof}
     Take a cover $\{U_\alpha\}_\alpha$ of $f^{-1}(D(h))$, and note that
     \begin{align*}
       f
     \end{align*}
   \end{proof}
 \end{proof}
 
 \begin{enumerate}
 	\item[3.3.] {
       $(*)$
       \begin{enumerate}
         \item Show that a morphism $f:X\to Y$ is of finite type if and only if it is locally of finite type and quasi-compact.
         \item Conclude from this that $f$ is of finite type if and only if for \emph{every} open affine subset $V=\Spec B$ of $Y$, $f^{-1}(V)$ can be covered by a finite number of open affines $U_j=\Spec A_j$, where each $A_j$ is a finitely generated $B$-algebra.
         \item Show also if $f$ is of finite type, then for \emph{every} open affine subset $U=\Spec A\subset f^{-1}(V)$, $A$ is a finitely generated $B$-algebra.
       \end{enumerate}
     }
 \end{enumerate}
 
 \begin{enumerate}
   \item[3.4.] {
       Show that a morphism $f:X\to Y$ is finite if and only if for \emph{every} open affine subset $V=\Spec B$ of $Y$, $f^{-1}(V)$ is affine, equal to $\Spec A$, where $A$ is a finite $B$-module.
     }
 \end{enumerate}
 
 \begin{proof}
   Fix an open affine cover $\{V_i\cong \Spec B_i\}$ of $Y$ such that $f^{-1}(V_i)\cong \Spec A_i$ is affine, and each $A_i$ is a finite $B_i$-module.
 \end{proof}
 
 \begin{enumerate}
   \item[3.6.] {
       Let $X$ be an integral scheme.
       Show that the local ring $\O_\xi$ of the generic point $\xi$ of $X$ is a field.
       It is called the \emph{function field} of $X$, and is denoted by $K(X)$.
       Show also if $U=\Spec A$ is any open affine subset of $X$, then $K(X)$ is isomorphic to the quotient field of $A$.
     }
 \end{enumerate}
 
 \begin{proof}
   Fix $U=\Spec A\subset X$ open affine.
   Then $\Cl_U(\{\xi\})=\{\xi\}^-\cap U=U$ implies that $\xi$ is the minimal prime ideal of $U$.
   That is, $\xi=(0)\in \Spec A$.
   Then $K(X)= A_{(0)}$ is the quotient field of $A$.
 \end{proof}

% \begin{enumerate}
%   \item[2.13.] {
%     }
% \end{enumerate}
% 
% \begin{proof}
% \end{proof}
% 
% \begin{enumerate}
%   \item[19.] {
%     }
% \end{enumerate}
% 
% \begin{proof}
% \end{proof}
% 
% \begin{enumerate}
%   \item[20.] {
%     }
% \end{enumerate}
% 
% \begin{proof}
% \end{proof}
% 
% \begin{enumerate}
%   \item[21.] {
%     }
% \end{enumerate}
% 
% \begin{proof}
% \end{proof}

\end{document}
